% ===== 原书 p.28 =====
\subsection{公理的验证}

由构造容易看出，a-ind 关于微分同胚是函子的，因此 a-ind 是指标函数。下面要说明解析指标满足两条公理：

\noindent\textbf{公理 A}\quad 若 $X$ 是一点，则 $\aind\pt:\Z\to\Z$ 是恒等映射。

\noindent\textbf{公理 B}\quad a-ind 与同态 $i_!$ 交换。

这两条公理唯一刻画拓扑指标 t-ind。于是只要它们成立，就得到 a-ind = t-ind，Atiyah–Singer 指标定理随即成立。

验证公理 A。若 $X$ 是一点，则 $X$ 上的椭圆算子只是一个线性变换 $P:V\to W$，其中 $V$、$W$ 是复向量空间，并且
\[
\sigma(P)=\dim E-\dim F\in K(TX)\cong\Z .
\]
此外注意
\[
\dim\Ker P+\dim\operatorname{Im}P=\dim E,\qquad
\dim F-\dim\operatorname{Im}P=\dim\Coker P .
\]
于是
\[
\Ind P=\dim\Ker P-\dim\Coker P=\sigma(P),
\]
公理 A 立即成立。

\transnote{这两条公理把指标定理拆成归一点与切除两步。要点在唯一性：任何满足两条公理的指标函数都等于 t-ind，t-ind 又按构造满足它们。所以证明 a-ind = t-ind 不必直接计算 a-ind，只需验证公理。}

验证公理 B 要麻烦一些。回忆紧子流形 $X\subset Y$ 的包含映射 $i:X\to Y$，我们曾把 $i_!$ 定义为合成
\[
i_!:K(TX)\xrightarrow{\ \varphi\ }K(TN)\xrightarrow{\ h_*\ }K(TY),
\]
其中 $N$ 是 $X$ 在 $Y$ 中的管状邻域，$\varphi$ 是 Thom 同构，$h_*$ 是由开包含 $h:TN\to TY$ 诱导的自然延拓同态。要验证 a-ind 与 $i_!$ 交换，只须证明它与 $\varphi$ 及 $h_*$ 都交换。先看自然延拓同态 $h_*$。

\begin{proposition}[切除]
若 $h:U\to X$ 是开集 $U$ 到 $X$ 的包含映射，则
\[
\aind\circ h_*(x)=h_*\circ\aind(x),
\]
其中 $h_*:K(TU)\to K(TX)$ 是自然延拓同态。
\end{proposition}

\begin{proof}
$K(TU)$ 中任一元素 $a$ 可由 $TU$ 上的齐次复形
\[
0\longrightarrow\pi^*E\xrightarrow{\ \alpha\ }\pi^*F\longrightarrow 0
\]
表示。$U$ 未必紧，故须取 $\alpha$ 为零次齐次的；还要求 $\alpha(x,\xi)=\mathrm{Id}$ 对所有位于某个紧集 $C\subset U$ 之外的 $x$ 成立（非紧情形的讨论见 [10, pp. 493–3]）。把 $E$、$F$ 平凡地延拓为 $X$ 上的丛 $E'$、$F'$，再考虑 $X$ 上的复形
\[
0\longrightarrow\pi^*E'\xrightarrow{\ \alpha\ }\pi^*F'\longrightarrow 0,
\]
% ===== 原书 p.29 =====
它表示元素 $h_*a\in K(TX)$，其中
\[
\alpha'=
\begin{cases}
\alpha, & \text{在 }TU\text{ 上},\\
\mathrm{Id}, & \text{在 }TX-TU\text{ 上}.
\end{cases}
\]
设 $P$、$P'$ 分别是 $U$ 与 $X$ 上对应的算子。由构造有 $\sigma(P')=h_*\sigma(P)$。若 $f$ 的支撑在 $U$ 内且 $Pf=0$，则 $f$ 的支撑必在 $U$ 内，从而 $P'f=0$。于是 $\Ker P=\Ker P'$。对伴算子 $P^*$ 与 $P'^*$ 同理，所以 $\Ind P=\Ind P'$。
\end{proof}

\transnote{原文取零次齐次并要求 $\alpha$ 在紧集外为恒等，是紧支撑 K 理论的约定：只有象征在无穷远处趋于恒等，才能把 $E$、$F$ 平凡延拓到 $X$ 而不改变 $K(TX)$ 中的类。若省掉紧支撑条件，延拓后的复形未必还表示 $h_*a$。}

剩下要证的是：解析指标与 Thom 同构 $\varphi:K(X)\to K(V)$ 交换，其中 $V$ 是 $X$ 上的复向量丛。回忆 $\varphi(x)=\pi^*x\cdot\lambda_V$，其中 $\pi:V\to X$ 是丛投影，$\lambda_V=\varphi(1)$ 是由外复形 $\Lambda(V)$ 给出的 Thom 类。要证的是
\[
\aind(\pi^*x\cdot\lambda_V)=\aind(x)\cdot\aind(\lambda_n)=\aind(x).
\]
为此须先建立 a-ind 的乘积公式，再在适当条件下验证 $\aind(\lambda_n)=1$。先从乘积公式入手。若丛可写成平凡丛 $V=X\times\R^n$，问题会大大简化。对一般向量丛，须改而考虑「扭」乘积。

设有主纤维丛 $P\to X$，其相伴的紧 Lie 群为 $H$：$H$ 从右自由作用在 $P$ 上，$X=P/H$。再设 $F$ 是局部紧空间，$H$ 从左作用在 $F$ 上，则得到 $X$ 上的纤维丛 $Y$：
\[
Y=P\times_H F=(P\times F)\big/\big((p,f)\sim(ph^{-1},hf)\ \ \forall h\in H\big).
\]
$H$ 在 $F$ 上的作用诱导 $TF$ 上的 $H$ 作用，因此可构造 $Y$ 上的向量丛 $P\times_H TF$。取度量后得到分解
\[
TY=(P\times_H TF)\oplus\pi^*TX,
\]
从而有乘积映射 $K(TX)\otimes K(P\times_H TF)\to K(TY)$。考虑同态
\[
K_H(TF)\to K_H(P\times TF)\cong K(P\times_H TF),
\]
其中 $K_H(\cdot)$ 是 §4 讨论的等变 K 理论函子，于是得到乘法映射
\[
K(TX)\otimes K_H(TF)\to K(TY).
\]
我们关心 $H=O(n)$、$F=\R^n$ 的情形。此时任一向量丛 $Y\to X$ 都可写成 $Y=P\times_{O(n)}\R^n$，其中 $P\to X$ 是适当的主纤维丛；于是得到乘法映射
\[
K(TX)\otimes K_{O(n)}(T\R^n)\to K(TY).
\]
下面给出解析指标的乘积公式。

\begin{proposition}[乘性公理]
若 $a\in K(TX)$，$b\in K_H(TF)$，则
\[
\aind^Y a\cdot b=\aind^H a\cdot\aind_H^F b,
\]
只要 $\aind_H^F b\in\Z$ 是平凡表示 $1\in R(H)$ 的整数倍。这里对 $\sigma(P)=b$ 定义 $\aind_H b=\Ind_H P=[\Ker P]-[\Coker P]\in R(H)$（见 §4）。
\end{proposition}

\transnote{等号左边省略了括号：$\aind^Y a\cdot b$ 指 $\aind^Y(a\cdot b)$。$a$ 是 $K(TX)$ 中的类，本身不在 $K(TY)$ 中，只有乘积 $a\cdot b$ 才有指标。}

% ===== 原书 p.30 =====
\begin{proof}
元素 $a$ 可由 $TX$ 上的 1 次齐次复形
\[
0\longrightarrow\pi^*E^0\xrightarrow{\ \alpha\ }\pi^*E^1\longrightarrow 0
\]
表示。构造 $X$ 上以 $\alpha$ 为象征的伪微分椭圆算子 $A$。取从属于平凡化 $Y$ 的开覆盖 $\{U_i\}$ 的单位分解，可把 $A$ 提升为 $Y$ 上的整体椭圆算子 $\tilde A$。\footnote{若 $P:\Gamma E\to\Gamma F$ 是 $X$ 上的算子，$G$ 是 $Y$ 上的向量丛，则 $P$ 到 $X\times Y$ 的提升（lift）$\tilde P:\Gamma(E\BT G)\to\Gamma(F\BT G)$ 定义为 $\tilde P(u(x)\otimes v(y))=P(u(x))\otimes v(y)$。}同理，$b$ 可由 $TF$ 上的 1 次齐次复形
\[
0\longrightarrow\pi^*G^0\xrightarrow{\ \beta\ }\pi^*G^1\longrightarrow 0
\]
表示。

设 $B$ 是 $F$ 上以 $\beta$ 为象征的椭圆伪微分算子，并要求 $B$ 与 $H$ 在 $G^0$、$G^1$ 上的作用交换。于是 $B$ 可提升为 $Y$ 上的整体椭圆算子 $\tilde B$。

现在构造 $Y$ 上的算子
\[
D=\begin{pmatrix}\tilde A & -\tilde B^*\\ \tilde B & \tilde A^*\end{pmatrix},
\qquad\text{其象征为}\qquad
\sigma(D)=\begin{pmatrix}\alpha\BT 1 & -1\BT\beta^*\\ 1\BT\beta & \alpha^*\BT 1\end{pmatrix}.
\]
回忆两个长度为 1 的复形之积的表达式，可知 $\sigma(D)$ 对应 $K(TY)$ 中的元素 $ab$，它由复形
\[
0\longrightarrow E^0\BT G^0\oplus E^1\BT G^1\xrightarrow{\ \sigma(D)\ }E^1\BT G^0\oplus E^0\BT G^1\longrightarrow 0
\]
表示。因此只须计算 $D$ 的指标。注意 $\tilde A$ 与 $\tilde B$ 交换，考察两个对角化算子
\[
D^*D=\begin{pmatrix}\tilde A^*\tilde A+\tilde B^*\tilde B & 0\\ 0 & \tilde A\tilde A^*+\tilde B\tilde B^*\end{pmatrix}
=\begin{pmatrix}P_0 & 0\\ 0 & Q_0\end{pmatrix},
\]
\[
DD^*=\begin{pmatrix}\tilde A\tilde A^*+\tilde B^*\tilde B & 0\\ 0 & \tilde A^*\tilde A+\tilde B\tilde B^*\end{pmatrix}
=\begin{pmatrix}P_1 & 0\\ 0 & Q_1\end{pmatrix},
\]
其中
\[
\Ker D=\Ker D^*D=\Ker P_0\oplus\Ker Q_0,
\]
\[
\Coker D=\Ker D^*=\Ker DD^*=\Ker P_1\oplus\Ker Q_1.
\]
于是
\[
\Ind D=(\dim\Ker P_0-\dim\Ker P_1)+(\dim\Ker Q_0-\dim\Ker Q_1).
\]
考察算子 $P_0=\tilde A^*\tilde A+\tilde B^*\tilde B$，注意
\[
\ip{P_0 u}{u}=\ip{\tilde A u}{\tilde A u}+\ip{\tilde B u}{\tilde B u},
\]
所以 $\Ker P_0=\Ker\tilde A\cap\Ker\tilde B$。$\tilde B$ 把 $B$ 延拓到 $Y$ 的各纤维，因此 $\Ker\tilde B$ 是 $X$ 上向量丛 $K_B=P\times_H\Ker B$ 的光滑截面空间。$\tilde A$ 诱导 $K_B$ 的截面上的算子 $C$，其象征 $\sigma(C)=\alpha\otimes\mathrm{Id}(K_B)$，于是 $\sigma(C)=a[K_B]\in K(TX)$，

% ===== 原书 p.31 =====
其中 $[K_B]$ 是 $K_B$ 在 $K(X)$ 中的类。把 $\tilde A$ 换成 $\tilde A^*$，$P_1$ 与 $C^*$ 有类似结论，于是
\[
\begin{aligned}
\dim\Ker P_0-\dim\Ker P_1
&=\dim\Ker C-\dim\Ker C^*\\
&=\Ind C=\aind^X(a[K_B]).
\end{aligned}
\]
同样，取 $L_B=P\times_H\Coker B$，得到
\[
\dim\Ker Q_0-\dim\Ker Q_1=\aind^X(a[L_B]).
\]
合并两式，得
\[
\Ind D=\aind^X\big(a\cdot([K_B]-[L_B])\big).
\]
若还有 $\aind_H^F b=[\Ker B]-[\Coker B]\in R(H)$ 是整数（即平凡表示 $\C$ 的整数倍），则 $[K_B]-[L_B]=(\aind_H^F b)\cdot 1$ 是平凡丛 $1$ 的整数倍。又因 a-ind$^X$ 是同态，即得所求的乘积公式
\[
\Ind^Y D=(\aind^X a)\cdot(\aind_H^F b).
\]
\end{proof}

\begin{proposition}[归一化公理]
若 $j:\pt\to\R^n$ 是原点的包含映射，于是有诱导同态 $j_!:R(O(n))\to K_{O(n)}(T\R^n)$，则
\[
\aind_{O(n)}j_!(1)=\aind_{O(n)}\lambda_n=1,
\]
其中 Bott 类 $\lambda_n=\lambda_{\C^n}$ 是 $\C^n$ 的 K 理论 Thom 类。
\end{proposition}

\begin{proof}
先看 $n=2$。回忆 de Rham 象征 $\rho_2=\rho_{S^2}$ 由 $2$-球面 $S^2$ 的 de Rham 复形复化得到。涉及的丛可定向，故结构群可取为 $SO(2)$。考察 $SO(2)$ 在 $S^2=\R^2\cup\infty$ 上的作用，可知 $\rho_2\in K_{SO(2)}(TS^2)$；经过适当形变得到（见 [10, §3.2]）
\[
\rho_2=h(\lambda_2)+f^*\circ h(\lambda_2),
\]
其中 $h$ 是延拓同态 $K(T\R^2)\to K(TS^2)$，$f:S^2\to S^2$ 是交换 $0$ 与 $\infty$ 的反射。a-ind 是函子的，且与自然延拓同态交换，故 $\aind_{SO(2)}\rho_2=2\aind_{SO(2)}\lambda_2$。注意 $\aind_{SO(2)}\lambda_2=\chi(S^2)=2$，于是 $\aind_{SO(2)}\lambda_2=1$。

\transnote{原文此处写「$\aind_{SO(2)}\lambda_2=\chi(S^2)=2$」，与上一句 $\aind_{SO(2)}\rho_2=2\aind_{SO(2)}\lambda_2$ 自相矛盾：$\lambda_2$ 是 Thom 类，$\chi(S^2)$ 应由 de Rham 象征 $\rho_2$ 给出。按上下文应为 $\aind_{SO(2)}\rho_2=\chi(S^2)=2$。译文照录原文，读者按此理解。}

$n=1$ 的情形。回忆我们讨论伪微分算子时提到的圆 $S^1$ 上的椭圆算子 $A$。经过适当同伦得到 $-h(\lambda_1)=\sigma(A)$，其中 $h$ 是自然延拓同态 $K(T\R^1)\to K(TS^1)$。由 $\Ind A=-1$ 可知 $\aind_{O(1)}\lambda_1=1$。这个构造在附录里作为粘合函数的例子详细展开。

把 $\R^n$ 分解成若干 $\R^1$ 与 $\R^2$ 的因子。由乘积公式可知，对 $O(n)$ 的每个子群 $G=G_1\times\cdots\times G_n$（各 $G_i$ 取 $O(1)$ 或 $SO(2)$）都有 $\aind_G\lambda_n=1$。$O(n)$ 的表示由它到这些子群的限制完全决定，故 $\aind_{O(n)}\lambda_n=1$。
\end{proof}

设 $\varphi:K(TX)\to K(TY)$ 是复向量丛 $TY\to TX$ 的 Thom 同构。把上面两个命题合起来得到
\[
\begin{aligned}
\aind^Y\circ\varphi(a)
&=\aind^Y(\pi^*a\cdot\lambda_{TY})\\
&=(\aind^X a)\cdot\big(\aind_{O(n)}^{\R^n}\lambda_n\big)=\aind^X a.
\end{aligned}
\]
至此已证解析指标与 Thom 同构交换，Atiyah–Singer 指标定理的证明完毕。
